\scanpage{34}
définie par l'\unsure{aut.|aut.t} \deleted{\unsure{de|$\tau_X$}} $\tau_X$ de $X(K)$). Ainsi, un
motif sur $\mathbb R$ définit :

a) Une structure de Hodge \inserted{$V$} sur $K$

b) Un autom. $f$\deleted{\unsure{pro}} \inserted{de carré 1} du $K$-espace vectoriel \deleted{sous-jacent}\inserted{$V_K$},
cet autom étant ``réel'' i.e. \deleted{\unsure{$f$|qui}} transformant $V_{\mathbb R}$ en $V_{\mathbb R}$,
\deleted{\unsure{On signifie}} et \unsure{$=$} \uline{rationnel} i.e. transforme $V$ en $V$

Ces deux données étant assujetties à la condition
\[
\left\{\begin{aligned}
 &( \deleted{f\unsure{\sigma}} )\,\sigma\ \text{inverse les filtrations naturelles}\\
 &\text{i.e. } f\sigma\ \text{échange la filtration donnée et l'opposée}
\end{aligned}\right.
\]

En termes de ``donnée de descente'', on peut \deleted{dire aussi}\inserted{exprimer}
aussi b) en disant qu'on a un isom
de la structure de Hodge $V$ avec la structure de
Hodge \deleted{\unsure{que}} ``conjuguée'' : c'est le frobénius réel.

Soit $G$ le groupe de Galois motivique de $M$ sur $\mathbb R$.
Alors $f$ correspond à un élément
\[
f_{\mathbb R}\in G(\mathbb Q)
\]
appelé \uline{frobénius réel}. C'est un élément d'ordre 2.
Son image dans $G/G^0$ engendre ce groupe. Cet
élément normalise $\underline i:S_{\mathbb R}\longrightarrow G_{\mathbb R}$, de façon précise
\[
f_{\mathbb R}\,i(\lambda)\,f_{\mathbb R}^{-1}=i(\deleted{\unsure{\sigma\lambda}}\lambda')
\qquad\text{où }\lambda\unsure{\sim|\mapsto}\lambda'\ \text{est la symétrie canonique de }S_{\mathbb R}
\]
donc sur $U_{\mathbb R}$ on a
\[
 f_{\mathbb R}\,i(\lambda)\,f_{\mathbb R}^{-1}=i(\lambda^{-1})
 \qquad(\lambda\in U_{\mathbb R}\deleted{\unsure{\mathbb R}}) \tag*{(*)}
\]

Le foncteur $M\longmapsto(\text{structure de Hodge associée, avec }f_{\mathbb R})$
est pleinement fidèle en vertu des conj. de Hodge.

% margin note written sideways
\margin{La structure en \deleted{\unsure{sont}} donne $G$ sur $\mathbb Q$, $\unsure{\varepsilon|E}:G\unsure{-|\to}G_m$, $i':S_{\mathbb R}\unsure{-|\to}G_{\mathbb R}$ (avec $i|G_{m\,\mathbb R}$ défini sur $\mathbb Q$ \unsure{etc|et}) $+\, f_{\mathbb R}\in G(\mathbb Q)$ satisfaisant la condition conj.}
% archive sheet number: 33
